\section{The marked order-four filling}\label{sec:order-four}

\subsection{The order-four datum and its conormal coordinate}

Consider the Weierstrass family over a disk about $t=1$,
\begin{equation}\label{of:eq:original}
 y^2=x^3-3t^3(t-1)x+2t^4(t-1)^2.
\end{equation}
The disk coordinate is $q=t-1$, with its complex orientation.
The origin is the section at infinity. Put
\begin{equation}\label{of:eq:reduction}
 q=s^4,\qquad x=s^2X,\qquad y=s^3Y.
\end{equation}
The reduced equation and discriminant are
\begin{equation}\label{of:eq:smooth}
 Y^2=X^3-3(1+s^4)^3X+2s^2(1+s^4)^4,
 \qquad \Delta=1728(1+s^4)^8.
\end{equation}
Thus its projective elliptic family is smooth on a sufficiently small disk.

The deck generator fixing $x,y,t$ is
\begin{equation}\label{of:eq:deck}
 (X,Y,s)\longmapsto(-X,iY,is).
\end{equation}
It multiplies $dX/Y$ by $i$. A counterclockwise meridian in the $q$ disk
lifts from $s$ to $is$. The endpoint identification uses the inverse deck
transformation, so its positive transport multiplier is $-i$.

We use the full integral period marking
\begin{equation}\label{of:eq:lattice}
 f_1=(1,0),\quad f_2=(i,0),\quad
 f_3=\tfrac12(1+i,1),\quad f_4=(0,\tau),\qquad \operatorname{Im}\tau>0.
\end{equation}
The positive multiplicative period and positive monodromy are
\begin{equation}\label{of:eq:T}
 \delta=-f_1-f_2+2f_3,
 \qquad
 T=\begin{pmatrix}0&1&0&0\\-1&0&-1&0\\0&0&1&0\\0&0&0&1\end{pmatrix}.
\end{equation}
Indeed multiplication by $-i$ on the first complex coordinate gives these
four columns and fixes the second coordinate. In particular $T^4=I$.

The ordered basis $(f_1,f_2,f_3,f_4)$ has the complex orientation.
We retain the logarithmic increment
\begin{equation}\label{of:eq:increment}
                         a_1=-f_4/4.
\end{equation}
This is the order-four increment paired with $a_0=e_4/3$ in
\cite[p.~12]{Engel}. The period marking is that of
\cite[Theorem~5.1, pp.~12--13]{Engel}.

The following hypotheses specify the geometric comparison.
Let $P_0$ and $O$ be disjoint sections near $t=1$, with $O$ the origin,
and let $\mathcal M=\mathcal O(P_0-O)$.
Choose a nonzero section $e$ of $\mathcal M|_O$ over this disk.
After \eqref{of:eq:reduction}, suppose the transformed sections $P'_0,O'$
are disjoint and put $\mathcal M'=\mathcal O(P'_0-O')$.

Assume the specified lifted translation on $\mathcal M^\times$ extends to
$(\mathcal M')^\times$ and commutes with the natural deck action.
Assume its quotient $\mathcal F$ is a proper smooth family of complex
two-tori over the $s$ disk.
Assume its smooth integral period marking extends \eqref{of:eq:lattice},
with multiplicative period $\delta$ and monodromy \eqref{of:eq:T}.

These are hypotheses on the reduction and its lifted translation.
They are not consequences of the matrix $T$.
The elliptic equation verifies smooth elliptic reduction directly.
The extension and properness of the specified lifted translation remain
separate hypotheses.
They are proved for the selected global construction in
Section~\ref{sec:holomorphic-realization}.

\begin{lemma}\label{of:lem:conormal}
There is a nonzero section $e'$ of $\mathcal M'|_{O'}$ such that the
canonical punctured divisor comparison gives
\begin{equation}\label{of:eq:conormal}
                     e=s^{-1}e',\qquad \rho e'(s)=-i\,e'(is).
\end{equation}
In the origin defined by $e'$, the original and twisted torus deck maps are
\begin{equation}\label{of:eq:affine}
 (z,s)\mapsto(T^{-1}z-\delta/4,is),\qquad
 (z,s)\mapsto(T^{-1}z-(\delta+f_4)/4,is),
\end{equation}
where torus coordinates are taken modulo their integral periods.
\end{lemma}

\begin{proof}
Near $O$, the divisor line $\mathcal O(P_0-O)$ is the ideal of $O$.
Its restriction to $O$ is therefore the conormal line.
The same statement holds at $O'$.

The local transverse parameters $w=-x/y$ and $W=-X/Y$ satisfy
$w=s^{-1}W$. These are parameters at infinity: the projective cubic in
the chart $Y\ne0$ solves $Z/Y$ holomorphically as a function of $X/Y$
near the origin. Write $e=u(q)w$ for a holomorphic unit $u$ and set
$e'=u(s^4)W$. This proves the first equality in \eqref{of:eq:conormal}.

Since $e$ is equivariant, the forward factor on $e'$ is
$(is)^{-1}/s^{-1}=-i$. Multiplication by $-i$ in the multiplicative fiber
is translation by $-\delta/4$.

An isomorphism between compact complex tori is affine in chosen origins.
Indeed its lifted derivative is lattice-periodic and descends to holomorphic
functions on a compact torus, hence is constant.
The period action therefore determines its linear part, here $T^{-1}$.
The image of $e'$ determines its translation by the preceding calculation.

The prescribed increment \eqref{of:eq:increment} gives the second formula.
The integral matrices are locally constant over the smooth disk.
The indicated translations are the specified torsion sections throughout it.
\end{proof}

This calculation uses the conormal character, not the normal-coordinate
character. It also retains the original section $e$.
On the positive root $s=e^{2\pi iv/4}$ that section contributes
$-v\delta/4$, and the logarithmic modification contributes $-vf_4/4$.
A nonunit boundary radius changes only a constant frame normalization.
The angular displacement is unchanged.

\subsection{The marked quotient and its collar}

Let $X=\mathbb R^3/\mathbb Z^3$ in the basis $(f_1,f_2,f_3)$, and write
$t_4\in\mathbb R/\mathbb Z$ for the fourth real period coordinate.
Set $U=(-1,-1,2)$ in $X$ and
\begin{equation}\label{of:eq:P}
 P=\begin{pmatrix}0&-1&-1\\1&0&0\\0&0&1\end{pmatrix},
 \qquad P^4=I,\qquad PU=U.
\end{equation}
Thus $T=P^{-1}\oplus1$.
The real disk below is identified with a closed complex disk by a positive
rescaling of $s$.

Define
\begin{equation}\label{of:eq:N}
 N=(X\times S^1_{t_4}\times D)/\langle g\rangle,
 \qquad g(x,t_4,s)=(Px-U/4,t_4-1/4,is).
\end{equation}
The action has order four and is free: for $1\le j\le3$, its $j$th power
changes $t_4$ by $-j/4$. The map $q=s^4$ descends to $N$.
Its reduced central fiber is the quotient at $s=0$.

Use the positive boundary mapping torus
\[
 B_T=(X\times S^1_{t_4}\times[0,1])/
 ((x,t_4,1)\sim(P^{-1}x,t_4,0)).
\]
The boundary identification determined by $e$ and \eqref{of:eq:increment} is
\begin{equation}\label{of:eq:j}
 j[x,t_4,v]=[x-vU/4,t_4-v/4,e^{2\pi iv/4}].
\end{equation}
At $v=1$, the inverse of $g$ sends this endpoint to
$[P^{-1}x,t_4,1]$. Hence \eqref{of:eq:j} has precisely the prescribed
positive return and starting fiber. Also $qj=e^{2\pi iv}$.

\begin{theorem}\label{of:thm:geometry}
Under the stated reduction hypotheses, \eqref{of:eq:N} is the marked
local filling for \eqref{of:eq:increment}.
Put
\[
 G_P=(X\times\mathbb R)/([y,\tau+1]=[Py,\tau]).
\]
There is an orientation-preserving diffeomorphism
\begin{equation}\label{of:eq:Phi}
 \Phi[x,t_4,s]=([x-t_4U,4t_4],e^{2\pi it_4}s)
                 \in G_P\times D.
\end{equation}
Write its second coordinate as $z$.
Then $q=e^{-2\pi i\tau}z^4$.

The fiber and positive boundary maps followed by retraction are
\begin{equation}\label{of:eq:fk}
 f(x,t_4)=[x-t_4U,4t_4],\qquad
 k[x,t_4,v]=[x-t_4U,4t_4-v].
\end{equation}
The integral product coordinates have fourth vector
\begin{equation}\label{of:eq:v4}
                         v_4=f_4+\delta.
\end{equation}
\end{theorem}

\begin{proof}
The integral shear $y=x-t_4U$ has determinant one.
It transforms $g$ into $(y,t_4,s)\mapsto(Py,t_4-1/4,is)$.
Both coordinates in \eqref{of:eq:Phi} are invariant under this map.
They are also invariant under $t_4\mapsto t_4+1$.

Given $([y,\tau],z)$, choose a real lift of $\tau$ and put
$t_4=\tau/4$, $x=y+t_4U$, $s=e^{-2\pi it_4}z$.
Changing the lift by one and replacing $y$ by $P^{-1}y$ changes the
result by $g^{-1}$. This proves a smooth inverse.
The equation for $q$ follows by taking fourth powers.

In these coordinates \eqref{of:eq:j} becomes
\[
 \Phi j[x,t_4,v]=([x-t_4U,4t_4-v],e^{2\pi it_4}).
\]
Retraction gives \eqref{of:eq:fk}. Holding $y$ fixed while increasing
$t_4$ gives the original vector $f_4+\delta$, proving \eqref{of:eq:v4}.

Orient $G_P$ by the $X$ orientation followed by increasing $\tau$.
Orient its disk by the complex orientation. The shear has determinant one,
the factor $4$ is positive, and the normal coordinate changes by a complex
rotation. Thus \eqref{of:eq:Phi} preserves the six-dimensional orientation.

The outward-first boundary orientation is the $X$ orientation followed by
$\tau$ and the positive argument of $z$.
The boundary circle matrix, in coordinates $(t_4,v)$, is
\begin{equation}\label{of:eq:boundary-matrix}
                   \begin{pmatrix}4&-1\\1&0\end{pmatrix}.
\end{equation}
Its determinant is one, proving the boundary orientation assertion.
\end{proof}

The change from $(f_1,f_2,f_3,v_4)$ to the original periods has matrix
\begin{equation}\label{of:eq:V}
 V=\begin{pmatrix}1&0&0&-1\\0&1&0&-1\\0&0&1&2\\0&0&0&1\end{pmatrix},
 \qquad \det V=1.
\end{equation}
Its inverse is the actual torus map $(x,t_4)\mapsto(x-t_4U,t_4)$.
Thus all subsequent cellular maps in the original marking are obtained
by composing their characteristic maps with this explicit diffeomorphism.

The disk bundle in \eqref{of:eq:Phi} has a specified smooth complex normal
trivialization. An explicit collar is
\[
 ([y,\tau],z,\epsilon)\longmapsto
 \Phi^{-1}([y,\tau],(1-\epsilon)z),\qquad |z|=1,\quad0\le\epsilon<1/2.
\]
The outward vector is $-\partial_\epsilon$.
The phase $e^{-2\pi i\tau}$ in $q$ is part of the marked fibration.
A smooth disk trivialization does not erase that phase.

Let $h$ be the positive core loop $[0,\tau]$ and let $\gamma$ be the
positive marked meridian. Formula \eqref{of:eq:fk} gives
\begin{equation}\label{of:eq:pi}
 k_\#(\gamma)=h^{-1},\qquad f_\#(v_4)=h^4,
 \qquad k_\#(\gamma)^4=f_\#(-f_4-\delta).
\end{equation}
These equalities concern based loops through the marked origin.
They preserve the counterclockwise base meridian even though its image
traverses the core negatively.

\subsection{Integral cellular chains with the actual monodromy}

The action on the three-dimensional period lattice is not a signed
permutation in the basis $(f_1,f_2,f_3)$.
We therefore construct cells on which the homeomorphism $P$ acts.
In $\mathbb R^3$, take the common refinement of the four tilings
$P^j(16^{-1}\mathbb Z^3+[0,1/16]^3)$, $0\le j<4$.
More explicitly, cut by every hyperplane
$(P^{-j}x)_a=m/16$, for $a=1,2,3$ and $m\in\mathbb Z$.
Its closed cells are convex rational polytopes.
They have diameter at most $\sqrt3/16$, because the original cubic tiling
is one of the four refinements.

Barycentric subdivision, followed by the
quotient by $\mathbb Z^3$, gives a finite triangulation $\mathcal K$ of $X$.
No cell is identified with itself along a nontrivial translation.
The construction is periodic, and $P$ permutes its polytopes and their
barycenters. Thus $P$ acts simplicially.

Let $D=C_*(\mathcal K;\mathbb Z)$ with its integral simplicial boundary $d$.
Orient each simplex once. On its generators $P$ sends each oriented simplex
to its oriented image, with the appropriate permutation sign.
Then $Pd=dP$ and $P^4=I$ exactly.
Set $D_j=0$ outside $0\le j\le3$ and put
\[
                         M=P^{-1},\qquad S=I+P+P^2+P^3.
\]
Here $P,M,S$ denote chain maps in every degree.

In suspended cells the interval precedes the simplex.
The fourth fiber coordinate is \eqref{of:eq:v4}, so the actual fiber complex is
\begin{equation}\label{of:eq:CQ}
 \begin{aligned}
 C_n&=D_n\oplus D_{n-1},&d_C(x,y)&=(dx,-dy),\\
 Q_n&=D_n\oplus D_{n-1},&d_Q(x,y)&=(dx+(P-I)y,-dy).
 \end{aligned}
\end{equation}
Here $C$ is the product CW complex for $X\times S^1$, and $Q$ is the
mapping-torus CW complex for $G_P$.
The fiber monodromy on $C$ is $M\oplus M$.

The positive boundary complex is
\begin{equation}\label{of:eq:B}
 \begin{split}
 B_n&=C_n\oplus C_{n-1},\\
 d_B(x,y,z,w)&=(dx+(M-I)z,-dy+(M-I)w,-dz,dw).
 \end{split}
\end{equation}
In particular these are geometric cellular complexes, with their actual
nonzero simplicial differential. No exterior lattice replaces $D$.

\begin{theorem}\label{of:thm:chains}
The fiber covering, positive boundary prism, and boundary retraction have
cellular representatives
\begin{equation}\label{of:eq:FHK}
 \begin{aligned}
 F_n(x,y)&=(x,Sy),\\
 H_n(x,y)&=(0,-Mx),\\
 K_n(x,y,z,w)&=(x,Sy-Mz).
 \end{aligned}
\end{equation}
They satisfy
\begin{equation}\label{of:eq:identities}
 d_QF=Fd_C,\qquad d_QH+Hd_C=F(M-I),\qquad d_QK=Kd_B.
\end{equation}
The formula for $H$ uses the positive base prism throughout.
\end{theorem}

\begin{proof}
In the coordinates $(y,t_4)$ the map $f$ is $[y,4t_4]$.
The $t_4$ interval crosses four core intervals.
After subdivision, their coefficients are $I,P,P^2,P^3$.
This proves the formula for $F$ on product cells.

The positive boundary homotopy is $[y,4t_4-v]$, with $v$ increasing
from zero to one. A simplex at $t_4=0$ travels from core level zero
to level minus one. In the standard core interval this is the negative
interval on $P^{-1}$ of that simplex. Its coefficient is $-M$.

For a suspended fiber cell, the two interval coordinates map to the one
coordinate $4t_4-v$. The image has one lower cellular dimension, so its
prism coefficient is zero. This proves the formula for $H$.
The positive boundary mapping torus combines $F$ and this homotopy, giving
$K(a,b)=F(a)+H(b)$.

Finally $(P-I)S=S(M-I)=0$ and $Pd=dP$.
Substitution in \eqref{of:eq:CQ} gives
$d_QH(x,y)+Hd_C(x,y)=((M-I)x,0)=F(M-I)(x,y)$.
Substitution in \eqref{of:eq:B} gives
\[
 d_QK(x,y,z,w)=(dx+(M-I)z,-Sdy+Mdz)=Kd_B(x,y,z,w).
\]
The identity for $F$ follows in the same way from $(P-I)S=0$.
\end{proof}

\begin{proposition}\label{of:prop:split}
There is an integral chain splitting $B\cong Q\oplus Q[1]$.
Here $Q[1]_n=Q_{n-1}$ with differential $-d_Q$.
A section of $K$ and an identification of its kernel are
\begin{equation}\label{of:eq:split}
 \begin{aligned}
 s_Q(x,b)&=(x,0,-Pb,0),\\
 \ker K_n&=\{(0,y,Sy,w)\},&
 \phi(0,y,Sy,w)&=(Py,w)\in Q[1]_n.
 \end{aligned}
\end{equation}
Consequently the relative cone
\[
 R_n=Q_n\oplus B_{n-1},\qquad d_R(q,b)=(d_Qq+Kb,-d_Bb)
\]
is integrally isomorphic to $\operatorname{Cone}(I_Q)\oplus Q[2]$.
\end{proposition}

\begin{proof}
One checks $Ks_Q=I$ and $d_Bs_Q=s_Qd_Q$ directly.
If $K(x,y,z,w)=0$, then $x=0$ and $z=PSy=Sy$.
The kernel differential in $(y,w)$ is
$(-dy+(M-I)w,dw)$.
Multiplying its first coordinate by $P$ gives
$(-d(Py)-(P-I)w,dw)$, precisely $-d_Q(Py,w)$.
All changes of coordinates have integral inverses.

The cone of the split projection is the displayed direct sum.
On $\operatorname{Cone}(I_Q)$, the homotopy $(u,v)\mapsto(0,u)$
contracts the identity cone. This also gives an explicit chain contraction
of the complementary relative summand.
\end{proof}

For completeness, the projection to $Q[2]$ and its inclusion are
\begin{equation}\label{of:eq:relative-maps}
 \begin{split}
 \Pi(q,(x,y,z,w))&=(Py,w),\\
 J(a,b)&=(0,(0,P^{-1}a,Sa,b)).
 \end{split}
\end{equation}
They commute with the stated differentials and satisfy $\Pi J=I$.
These algebraic coordinates do not fix the sign of a geometric Thom class
without a comparison of orientations.

\subsection{Relative singular chains and the positive Thom class}

Here we give that geometric construction directly.
Use normalized singular cubical chains with integer coefficients.
A generator is a continuous map $[0,1]^n\to Z$, modulo cubes independent
of a coordinate. The boundary is
$\sum_{j=1}^n(-1)^{j-1}(\partial_j^1-\partial_j^0)$.

Write $\mathscr Q=C_*^\square(G_P)$ and
$\mathscr B=C_*^\square(B_T)$.
Let $i=j:B_T\to N$ be the boundary inclusion, let $r:N\to G_P$ be
retraction, and let $c:G_P\to N$ be the zero section.
The actual boundary chain map is $\mathscr K=(ri)_*$.
These maps specify the relative singular complex of the same filling.

Let $\eta$ be the positive normal-circle cube
$u\mapsto e^{2\pi iu}$.
Through \eqref{of:eq:Phi}, identify the boundary with $S^1_z\times G_P$,
putting the normal circle first for this chain notation.
Define
\begin{equation}\label{of:eq:L}
 L(q)=(\Phi i)^{-1}_*(\eta\times q).
\end{equation}
The inverse includes the interchange of the two product factors.
Then $dL=-Ld$, and $\mathscr K L=0$ exactly: projection makes the first
cubical coordinate degenerate.

Define the radial homotopy on a boundary cube $b$ by
\[
 \mathcal R(b)(u,\mathbf t)
   =\Phi^{-1}(r i b(\mathbf t),u\,z(i b(\mathbf t))).
\]
The radial interval is first. Its boundary identity is
\begin{equation}\label{of:eq:radial}
                 d\mathcal R+\mathcal R d=i_*-c_*\mathscr K.
\end{equation}

Therefore
\begin{equation}\label{of:eq:cone-comparison}
 (q,b)\longmapsto(c_*q-\mathcal R b,b)
 \colon\operatorname{Cone}(\mathscr K)\longrightarrow
                         \operatorname{Cone}(i_*)
\end{equation}
is a chain map. Its projection to $C_*^\square(N,B_T)$ is a
quasi-isomorphism. Indeed retraction and the zero section are inverse
homotopy equivalences, while $\operatorname{Cone}(i_*)$ maps to the
relative quotient by $(a,b)\mapsto[a]$. The latter has an acyclic kernel,
the identity cone on $C_*^\square(B_T)$.

\begin{proposition}\label{of:prop:Thom}
The positive Thom map is represented in the boundary cone by
\begin{equation}\label{of:eq:Thom}
 \Theta(q)=(0,-Lq)\in\operatorname{Cone}(\mathscr K),
 \qquad \deg\Theta=2.
\end{equation}
Its image under \eqref{of:eq:cone-comparison} and relative projection is
$[D_z]\times q=q\times[D_z]$, with the positive complex disk orientation.
It induces $H_{n-2}(G_P;\mathbb Z)\cong H_n(N,B_T;\mathbb Z)$.
\end{proposition}

\begin{proof}
Since $dL=-Ld$ and $\mathscr K L=0$,
$d(0,-Lq)=(0,dLq)=(0,-Ldq)$.
Thus \eqref{of:eq:Thom} is a degree-two chain map.
Its relative image is $\mathcal R Lq$.
The first two coordinates are
$(u,v)\mapsto u e^{2\pi iv}$, the positive disk cube.
Its only nondegenerate boundary is the positive normal circle.
The remaining coordinates are exactly those of $q$.

Interchanging a two-dimensional disk and a $k$-cube has sign $(-1)^{2k}=1$.
This proves the orientation statement in every degree.

The relative cubical cross product with the generator of
$H_2(D,S^1;\mathbb Z)$ is an isomorphism by the relative product
chain equivalence. Since this group is free and concentrated in degree two,
there are no additional tensor or torsion summands.
\end{proof}

The plus sign on a normal-circle insertion into the same cone would give
the negative Thom class. In particular the sign in
\eqref{of:eq:Thom} follows from the cone convention, not from a reversal
of the positive base meridian.
Cellular characteristic maps and cubical subdivision compare the cellular
and singular complexes of $N$ and $B_T$.
The constructions above give both the finite cellular relative cone and
the actual geometric Thom map. No zero-differential model is needed.

\subsection{Integral consequences in the original periods}

The action of $P$ on $H_*(X;\mathbb Z)$ is its exterior action.
This assertion is used here only for homology, after the geometric chain
maps have been constructed. In lexicographic exterior bases its matrices
in degrees one and two, and the corresponding norm maps, are
\begin{equation}\label{of:eq:exterior}
 \begin{aligned}
 P_1&=\begin{pmatrix}0&-1&-1\\1&0&0\\0&0&1\end{pmatrix},&
 S_1&=\begin{pmatrix}0&0&-2\\0&0&-2\\0&0&4\end{pmatrix},\\
 P_2&=\begin{pmatrix}1&1&0\\0&0&-1\\0&1&0\end{pmatrix},&
 S_2&=\begin{pmatrix}4&2&-2\\0&0&0\\0&0&0\end{pmatrix}.
 \end{aligned}
\end{equation}
In degrees zero and three, $P=1$ and $S=4$.

Both $P_1-I$ and $P_2-I$ have Smith factors $1,1$.
Their cokernels and kernels are therefore free of rank one.
The kernel generators are $\delta$ and $f_1\wedge f_2$, respectively.
The cokernel generators may be taken as $[f_3]$ and
$[f_1\wedge f_3]$.
In the second cokernel,
$[f_1\wedge f_2]=2[f_1\wedge f_3]$ and
$[f_2\wedge f_3]=-[f_1\wedge f_3]$.

\begin{proposition}\label{of:prop:homology}
The groups $H_n(N;\mathbb Z)$ are free of ranks $1,2,2,2,1$ in degrees
zero through four, and vanish otherwise.
The boundary groups are free of ranks $1,3,4,4,3,1$.
The map $k_*$ is a split surjection in every degree.
The covering $f_*$ has cokernels
\begin{equation}\label{of:eq:indices}
 0,\quad\mathbb Z/4,\quad\mathbb Z/2,\quad\mathbb Z/2,
 \quad\mathbb Z/4
\end{equation}
in degrees zero through four.
Its kernel ranks in those degrees are $0,2,4,2,0$.
The cohomological specialization indices in the full invariant lattices
are $1,4,2,2,4$ in degrees zero through four.
\end{proposition}

\begin{proof}
The mapping-torus exact sequence gives
\[
 0\to\operatorname{coker}(P_n-I)\to H_n(G_P;\mathbb Z)
   \to\ker(P_{n-1}-I)\to0.
\]
The groups at both ends are free by \eqref{of:eq:exterior}, including
degrees zero and three. Hence these extensions split as abelian groups,
and the asserted ranks follow. The boundary statement and split
surjectivity follow from Proposition~\ref{of:prop:split}.

The covering respects this exact sequence. On the first summand it induces
the quotient from $H_n(X)$ onto $\operatorname{coker}(P_n-I)$, which is
surjective. On the quotient term it induces the norm
$S:H_{n-1}(X)\to\ker(P_{n-1}-I)$.
Consequently its cokernel equals the cokernel of that norm, independently
of any choice of splitting.

The images of the four norm maps are
$4\mathbb Z$, $2\mathbb Z\delta$, $2\mathbb Z(f_1\wedge f_2)$,
and $4\mathbb Z$.
This proves \eqref{of:eq:indices}. Rank subtraction proves the kernel ranks.

The groups are free. The kernel of $f_*$ is saturated because its target
is free, and it equals the saturation of $(T-I)H_n(F)$ over $\mathbb Q$.
Thus its integral annihilator is exactly $H^n(F;\mathbb Z)^T$.
Dualizing a Smith normal form of $f_*$ gives the stated specialization
indices in that annihilator.
\end{proof}

In degree four the covering has geometric degree $+4$.
If $\omega$ is the oriented three-torus class, the interval-first core
class $\sigma_\tau\omega$ is the negative of its fiber-first orientation.
Consequently
\[
                  f_*(f_1\wedge f_2\wedge f_3\wedge f_4)
                        =-4\sigma_\tau\omega.
\]

The induced maps have a particularly simple deciding expression in degree one.
Let $c=[f_3]$ in the core and retain the positive core class $h$.
Then
\begin{equation}\label{of:eq:H1}
 f_*(f_1)=f_*(f_2)=0,\quad f_*(f_3)=c,\quad
 f_*(f_4)=-2c+4h,\quad k_*(\gamma)=-h.
\end{equation}
Indeed $f_4=v_4-\delta$ and $[\delta]=2c$.
Since $c$ has infinite order, the local relation
$\gamma^4=f_\#(-f_4)$ would contradict \eqref{of:eq:H1}.
The correct relation is \eqref{of:eq:pi}.

For comparison with the common period marking, let the columns of
\[
 C=\begin{pmatrix}1&1&0&1\\0&1&0&0\\0&-1&1&0\\0&0&0&1\end{pmatrix}
\]
express $(f_1,f_2,f_3,f_4)$ in $(e_1,e_2,e_3,e_4)$.
Then $C\delta=\delta_0=-2e_1-e_2+3e_3$ and
\begin{equation}\label{of:eq:common-vector}
 C(-f_4-\delta)=e_1+e_2-3e_3-e_4.
\end{equation}
This is an integral change of the actual period coordinates.
The invariant covector $\psi=e_4^*$ evaluates to $-1$ on this vector,
the same value as on $-f_4$. Equality of that scalar does not identify
the two marked local relations.

The order-four construction uses its stated smooth reduction, commuting
lift, and period marking. It constructs their marked filling and all its
boundary and relative chain maps. Section~\ref{sec:common-chain}
constructs the comparison maps and homotopies for both finite
monodromies. The cusp neighborhood is supplied separately in
Section~\ref{sec:toric-cusp}. Its boundary map, together with these
finite maps, determines the global integral gluing calculation.
